\documentclass[10pt]{article}
\usepackage[margin=0.85in]{geometry}
\usepackage{amsmath,amssymb,booktabs,float,graphicx,microtype,array,longtable,url,xcolor}
\microtypesetup{expansion=false}
\usepackage{xurl}
\emergencystretch=2em
\newcolumntype{L}[1]{>{\raggedright\arraybackslash}p{#1}}
\usepackage[colorlinks=true,allcolors=blue!45!black,breaklinks=true]{hyperref}
\usepackage[numbers,sort&compress]{natbib}
\hypersetup{pdftitle={Exact Filtered Vector Search Across Execution Environments},pdfauthor={Akshat Singh Kushwaha}}
\makeatletter
\let\tableinput\@@input
\makeatother
\newcommand{\qenlo}{Qenlo}
\newcommand{\E}{E}
\newcommand{\D}{D}
\newcommand{\B}{B}
\newcommand{\fig}[3]{\begin{figure}[H]\centering\includegraphics[width=\linewidth]{figures/final/#1.pdf}\caption{#2}\label{#3}\end{figure}}
\title{Exact Filtered Vector Search Across Execution Environments:\\Why Scalar Routing Rules Fail}
\author{Akshat Singh Kushwaha\\\texttt{akshatsingh14372@outlook.com}\\\url{https://github.com/a3ro-dev/qenlo}}
\date{September 2026}
\begin{document}
\maketitle
\begin{abstract}
Filtered vector search chooses an execution route after a predicate restricts the candidates. We study that decision in Qenlo, an embedded vector store with exhaustive CPU, portable WGPU, and optional tensor execution. An archive-wide reduction inspects 303 sample files or archive members, leaving 282 content-distinct series and 1,826,130 CSV sample rows across twelve schemas; these counts describe coverage, not independent trials. A scalar $EDB$ rule selected with zero errors on 31 development CPU/GPU pairs fails a preregistered RTX 4050 suite on 8 of 16 workloads and reaches 235.7\% maximum backend-choice regret. Near-equal $EDB$ values with equal $k$ still produce different winners. Separate 384-dimensional, batch-one cohorts also disagree at the same eligible counts: RTX 4050 CPU/WGPU reverses between 2K and 3K rows, whereas an archived RTX 4090 cohort reverses between 4K and 6K. A certified FP32 CPU path improves A40 P95 by 17.0\% and 28.0\% at 1K and 4K eligible rows but regresses 10.0\% at 100K. One-pass eligibility preparation, corrected 768-dimensional results, lifecycle measurements, ANN comparisons, failures, and mobile aggregates supply mechanism and boundary evidence. The conclusion is negative but scoped: the frozen $EDB$ rule fails its gate, and a universal $E$ threshold cannot classify the observed winners. A post-hoc $EDB$ threshold has only 3.2\% maximum regret, so the archive does not reject every scalar policy at practical regret limits. Query shape, $k$, device/source identity, preparation, representation, and state motivate a new end-to-end held-out study; no successful router is yet validated.
\end{abstract}

\section{Introduction and research questions}
A filtered query over 100,000 vectors may score only one hundred candidates. The remaining work can be dominated by finding those candidates, dispatching a device, selecting results, or returning them through an application interface. Conversely, an unfiltered batch may amortize dispatch enough to favor exhaustive GPU execution even for a modest collection. Collection size and GPU availability alone provide an incomplete description of the execution problem.

Qenlo supplies a concrete system in which to examine this boundary. Canonical records are shared by rebuildable search structures; changing routes should preserve the logical result contract. The archive contains several generations of implementation and experiment design. Earlier studies emphasize eligible cardinality, predicate representation, and routing regret. The September 5 campaign adds small collections, tensor adapters, resource diagnostics, lifecycle measurements, and a rejected selector. A later preregistered router suite supplies an actual held-out falsification. This paper recursively inventories retained sample series and reduces previously contextual Phase 0 and Phase 2 archives without treating historical or candidate results as measurements of the final source tree.

We ask three questions. Under what measured conditions does the preferred exhaustive route change? How much of a completed call comes from eligibility preparation, selection, synchronization, or derived-state preparation rather than scoring? What deployment conclusions survive differences in numerical accuracy, APIs, persistence, runtime ownership, and evidence quality?

The contribution is a systems study, not a new asymptotic nearest-neighbor algorithm. It combines localized crossovers, an eligibility-materialization ablation, a failure-preserving small-collection campaign, heterogeneous-device observations, a sparse-only CPU optimization, and development/held-out routing evaluation. A claim ledger records which measurements support each conclusion. Rejected policies and optimizations retain scientific value, while their latency does not become a promise about the retained implementation.

\section{System model and semantic contracts}
Let the live canonical collection at generation $g$ contain normalized FP32 vectors $x_i\in\mathbb{R}^D$, public identifiers $i$, and metadata $m_i$. A predicate $P$ defines
\begin{equation}
 S_P(g)=\{i:\ i\text{ is live at }g\ \land\ P(m_i)\},\qquad E=|S_P(g)|.
\end{equation}
For query $q$, exhaustive search evaluates eligible rows and returns up to $k$ results ordered within each implementation by computed cosine distance and identifier. The ideal cosine distance is
\begin{equation}
 d(q,x_i)=1-\frac{q^\top x_i}{\|q\|_2\|x_i\|_2}.
\end{equation}
A batch contains $B$ queries. Scoring work scales with $BED$, but that expression omits predicate traversal, memory access, selection, and completed-call overhead.

\paragraph{Exhaustiveness and numerical agreement.}
We use ``exact'' in the algorithmic sense of exhaustive candidate coverage. Stored FP32 values, CPU accumulation, GPU arithmetic, and an FP64 reference need not yield identical rankings near ties. Oracle recall is reported independently. Historical dense native rows have recall@10 of 0.99998 despite exhaustive execution; corrected small-campaign rows have recall 1.0. Neither proves universal FP64 equivalence. ANN remains approximate even when observed recall is one. Deterministic engine ordering does not prove common FP64-oracle tie membership.

\paragraph{Canonical and derived state.}
The Rust store owns vectors, identifiers, metadata, tombstones, and generation. CPU scan layouts, WGPU residency, ANN indexes, and eligibility plans are derived state. Canonical vectors are FP32, normalized with an FP64 norm accumulator at ingestion and query boundaries. Restoration validates stored bytes without renormalizing them. Durable collections add snapshot and write-ahead-log behavior; in-memory benchmark collections have no durable generation. Source inspection supports these mechanisms, not universal power-loss durability or arbitrary concurrent or crash schedules.

\section{Execution paths}
\paragraph{CPU and portable GPU.}
Historical CPU artifacts describe FP64 accumulation with AVX2/FMA where available. Audited current CPU code also uses a bounded FP32 candidate stage with FP64 rescoring and reference fallback. Not every CPU arithmetic operation is FP64. These are Qenlo implementations, not bounds on optimized CPU performance. WGPU uses WGSL shaders and native graphics backends. WebGPU and WGSL specify a programming model, not performance portability \citep{webgpu,wgsl}; positive execution evidence is limited to named devices.

Resident GPU vectors use compact eligible rows, dense masks, or supported shader predicates. Bounded chunks emit candidates for host merging. Under a common total order, an item outside its chunk's top $k$ cannot enter the global top $k$, since at least $k$ items in that chunk precede it. This standard reduction explains bounded readback; it does not remove floating-point differences.

\paragraph{Eligibility compilation.}
An \texttt{EligibilityPlan} binds a predicate result to generation and records eligible count, representation, row IDs, and preparation diagnostics. One-pass compilation avoids counting eligibility and traversing it again to materialize rows. A bounded cache keyed by generation and predicate reuses plans only while that identity remains valid. Exact row execution validates ordering, bounds, liveness, and predicate membership. The historical ablation holds GPU scoring fixed across two-pass, one-pass, and warmed cached paths; it does not recalibrate the earlier laptop's crossover.

\paragraph{Tensor execution.}
The optional Python \texttt{TorchIndex} owns vector and ID tensors and supports CPU, CUDA, and MPS selection in source; only CPU and CUDA have measurements in this paper. The repository's conformance test is CPU-focused, so source-level device support should not be read as positive evidence for every CUDA or MPS device. A collection-derived snapshot records its generation and rejects staleness at search entry during sequential use. It neither holds a native collection lock throughout computation nor rechecks afterward; concurrent mutation guarantees are unestablished. A detached array index has no durable collection identity. Both are derived computation, not another durable database. Replays prepare filtered matrices before query timing. Tensor allocation fields exclude allocator and library overhead.

\paragraph{Mutation and route observability.}
Resident updates append rows or update liveness without rebuilding when capacity and layout permit; rebuild remains the fallback. CPU-only operation avoids GPU initialization. Required-GPU mode returns errors on failure, while automatic mode records fallback. Static routing and hardware-bound profiles exist, but no retained experiment validates a calibrated profile on held-out decisions. File-backed commits validate and stage durability before publication; this statement does not apply to in-memory collections.

\section{Methodology and evidence boundaries}
A comparison needs source and archive identity, hardware and software, dataset and query partition, $N,D,B,k,E/N$, predicate representation and semantics, API, timing, and correctness gate. Numerical observations stay within named cohorts; cross-cohort conclusions are qualitative synthesis. Historical endpoints with differing revisions remain descriptive context, not controlled ablations.

\paragraph{Timing.}
Native timing starts at the host search call and ends with host-visible ordered results, including in-call eligibility, transfers, dispatch, synchronization, readback, and merge. Load, construction, warmup, and oracle work are excluded. External small-campaign adapters also return host-visible identifiers and distances but receive a prefiltered canonical matrix. Their build field includes matrix and index construction; the metadata traversal creating eligible IDs precedes that timer. Even \texttt{build\_and\_filter\_prepare\_ns} omits part of preparation. Chroma includes Python collection bindings and serialization. These are scoped completed calls, not identical database APIs.

CUDA events are separate from host calls. Scoring, selection, and backend diagnostics overlap: summing them or subtracting them from P95 invents an unmeasured phase. Lifecycle records separate mutation, collection opening, and first search. Throughput fields including validation in their wall interval are not inverse query latency and are not used for cross-engine QPS rankings.

\paragraph{Samples and correctness.}
New synthetic cells generally use three repetitions of 64 evaluation queries: 192 batch-one calls, 24 batch-eight calls, or three batch-64 calls. The real fixture uses 5,000 queries per repetition and 939 batch-16 calls across three runs. Matrix \texttt{sample\_count} means native batch calls but external replay repetitions; Table~\ref{tab:headline} restores common calls and runs notation. Per-run nearest-rank percentiles are reduced by the lower median. Replays preserve native query order and truth and validate membership, distances, eligibility, and recall outside timing. Repeated queries are not independent deployment workloads.

Historical Windows and eligibility cells use five runs of 5,000 held-out calls, $k=10$, $B=1$, and 200 warmups; the statistic is median run P95. A6000 tables summarize per-call distributions within cells; synthetic A6000 has 500 calls per system and cell. Android and Arc provide aggregate percentiles. No figure claims significance or draws error bars. Lines are guides. Sequential execution, post-hoc choices, and small run counts limit causal and statistical interpretation.

The archive-wide census hashes canonical-LF CSV bytes. It includes tracked and selected archive-internal \texttt{samples.csv}, \texttt{raw\_samples.csv}, and lifecycle records; exact mirrors contribute one unique series. The twelve schemas have different sample units. Their 1,826,130 rows are never pooled into a confidence interval or treated as independent query draws. Numerical synthesis below uses only source-compatible comparisons, while the complete disposition ledger remains machine readable.

\section{Evidence cohorts and source identity}\label{sec:cohorts}
\begin{table}[H]\centering\small
\caption{Cohort map. Observations describe archive volume, not independent replications. Codes identify source roles below.}
\begin{tabular}{L{.24\linewidth}L{.23\linewidth}L{.22\linewidth}L{.19\linewidth}}
\toprule Cohort & Hardware / data & Scope & Evidence \\
\midrule
Small S0/S1/S2 & Ada, A40, 4090; generated, AG News & 1K--100K, 128--768D & 182 status rows \\
Localization H1 & RTX 4050/DX12; AG News & 100K, 384D, $E=2$K--10K & 300,000 raw calls \\
Endpoints H0 & RTX 4050; AG News & 100K, sparse/dense & 225,000 raw calls \\
Eligibility H2 & RTX 4090/Vulkan; AG News & 100K, 384D, three modes & 300,000 raw calls \\
Linux H3 & Cloud CPU; AG News & 100K, 384D, seven engines & 10,500 raw calls \\
A6000 real H4 & CUDA; TopK ModernBERT & 1M, 768D, strict predicate & 20,000 raw calls \\
A6000 synthetic H5 & CUDA; generated FP32 & 1M, 768D, four $E$ & 4,000 raw calls \\
Android H6 & Mali-G615/Vulkan & 10K/100K, 384D & 21 aggregate cells \\
Intel Arc H7 & Arc/Vulkan & 10K/100K, 384D & 21 aggregate cells \\
Router R1 & RTX 4050/DX12; generated & 16 held-out shapes & 48 raw route series \\
Phase 0 P0 & RTX 4090/Vulkan; AG News & 100K, 384D, $E=1$K--100K & 19 cells \\
Phase 2 P2 & A40 CPU; AG News & 100K, 384D, three $E$ & 8 reduced rows \\
\bottomrule\end{tabular}\end{table}
S0 is the frozen archive \texttt{f37e744c}, S1 is the experimental archive \texttt{730c0500}, and S2 is the corrected supplement \texttt{bb195fed}. Both S1 and S2 contain the lane-minimum candidate. The local WGSL is byte-identical to S0's simpler selector. The manifests describe a dirty worktree around HEAD \texttt{b22d3b5}; the archive hashes identify the measured source more precisely than that commit alone. The matrix label \texttt{current-gpu} is an artifact key, not a claim about final local behavior.

H1 records common revision \texttt{3e2a4a9} in a dirty worktree; its configurations and outputs are part of the evidence identity. H0 contains several revisions. H2 is archive \texttt{26e0ccc0}, and H3 is \texttt{bb51987}, mirrored and counted once. H4 and H5 use different A6000 environments and research CUDA implementations. H6 and H7 identify device-lab run IDs rather than reconstructible executable archives. R1 is SHA-256 \texttt{8e3b5fe6}; P0 is \texttt{7292fd1c}; and P2 is \texttt{ac3afda4}. Their executable revisions and hosts differ, so cross-cohort synthesis uses winner contradictions rather than pooled latency.

Windows data derive from DTU AG News embeddings \citep{nielsen2022agnews}; Linux uses a separately prepared fixture. Equal dataset names do not imply equal bytes or partitions. The small campaign retains its own fixture and checksums. The A6000 cohort uses TopK Bench's million-vector, 768-dimensional ModernBERT corpus \citep{topkbench}; its strict and provider-compatible predicates are separate. Appendix~\ref{sec:artifact} expands the provenance and exclusions.
\section{Small-collection campaign}
\subsection{Coverage and common sweep}
The matrix has 182 rows: 131 completed, 42 marked \texttt{failed\_or\_unavailable}, seven failed, and two marked invalid-harness. Only 130 rows qualify. The completed USearch batch-64 row remains visible at recall 0.3453125. The manifest records fifteen configurations: ten primary offers and five reference, deep, retry, or fix follow-ups. These are neither fifteen device replications nor ten successful devices. Results came from RTX 2000 Ada, two A40 configurations, and a reference RTX 4090; unavailable L4, RTX 3090, A6000, and community offers remain visible.
\begin{table}[H]\centering\small
\caption{Common synthetic workloads, $k=10$. Six selected cells, not a factorial sweep.}
\begin{tabular}{lrrrr}\toprule Cell & $N$ & $D$ & $B$ & $E/N$ \\\midrule
\tableinput{tables/final/common_definitions}
\bottomrule\end{tabular}\end{table}
\fig{small_common}{S1 \texttt{730c0500}, synthetic common cells, native completed-call P95. Panels are separate hosts. WGPU uses the candidate; no interpolation/error bars. Current-CPU rows are absent on A40, so those panels do not establish CPU/GPU pairs. All displayed FP64 recalls are 1.0. Vertical axis is logarithmic.}{fig:small-common}
On the reference RTX 4090, CPU wins C1, C2, C4, and C5, while WGPU wins the dense batched cells C3 and C6. In C5, 100K total rows leave 1K eligible: CPU takes 0.0341 ms P95 versus WGPU's 0.0953 ms. In C6, WGPU takes 4.0479 ms versus CPU's 68.0575 ms. Dimension, batching, and eligibility change together, so these are route observations rather than isolated size effects. The A40 configurations should not be combined into a pooled device curve.

\subsection{Corrected 100K-by-768 supplement}
S2 ran on Linux with an RTX 4090 and Vulkan, driver 580.159.04, and an AMD Ryzen Threadripper 7960X. The generated data record has CRC32 \texttt{50b69280}. Table~\ref{tab:headline} verifies all twelve rows against the raw records and matrix. These rows used the candidate subsequently rejected by the local source.
\begin{table}[H]\centering\small
\caption{S2, 100K$\times$768, completed-call ms. FP64 recall 1.0 throughout. Calls/runs counts batch calls across runs; replays follow native batches. External matrices are prefiltered outside timing.}\label{tab:headline}
\begin{tabular}{llrrrr}\toprule Workload & Engine & P50 & P95 & Recall & Calls/runs \\\midrule
\tableinput{tables/final/headline}
\bottomrule\end{tabular}\end{table}
\fig{small_768}{S2 \texttt{bb195fed}, RTX 4090/Vulkan, generated 100K$\times$768. Host-call P95, logarithmic horizontal axes, no error bars. Eligibility/batching differ between panels. WGPU is the rejected candidate; all FP64 recalls 1.0. FAISS Flat, NumPy, Torch use prefiltered matrices.}{fig:small-768}
For the unfiltered batch-one workload, WGPU is 10.49 times faster than native CPU at P95, while Torch CUDA is 1.34 times faster than WGPU. At 10\% eligibility and batch eight, the corresponding ratios are 31.91 and 2.33. These are archived calls. Fresh matched runs are required before attaching the WGPU values to the final full-rescan selector. Tensor latency excludes deployment and persistence costs around a matrix.

\subsection{Real embeddings}
The S1 real AG News workload has $N=100$K, $D=384$, $B=16$, $k=10$, and $E/N=1\%$. WGPU P95 is 0.171958 ms, CPU is 0.941633 ms, FAISS Flat is 0.242735 ms, NumPy is 0.749214 ms, Torch CPU is 0.268323 ms, and Torch CUDA is 0.351387 ms; all have FP64 recall 1.0. WGPU leads Torch CUDA in this workload, reversing their order in distinct S2 synthetic cells. The change across workload and source conditions does not isolate its cause. One fixture cannot establish robustness to different models, geometries, or correlated metadata. ANN rows and their recall remain in Appendix~\ref{sec:tables}.

\section{Historical crossover and eligibility preparation}
\subsection{A localized exhaustive winner reversal}
Matched H1 searches use 100K$\times$384 AG News on an RTX 4050 with DX12, $B=1$, and $k=10$. CPU takes 0.5198 ms P95 versus compact-row WGPU's 0.6340 ms at $E=2$K. At 3K, WGPU takes 0.8144 ms versus CPU's 0.9563 ms and remains faster through 10K. The reversal is bracketed by $(2\mathrm{K},3\mathrm{K})$; it is not a fitted threshold. All localization recalls are 1.0. Compact-row timing includes the duplicated count and materialization traversal.
\fig{phase_map}{Separate historical panels. Left: H1 \texttt{3e2a4a9}, RTX 4050/DX12, AG News 100K$\times$384, $B=1,k=10$, compact rows, median five-run P95, FP64 recall 1.0. Older endpoint revisions omitted. Right: H5 A6000 research CUDA/FAISS, generated 1M$\times$768, prefiltered rows, $B=1,k=10$, P95 over 500 calls/system/cell; unordered top-10 agreement without independent FP64 oracle. Logarithmic axes; lines guide the eye, not fitted laws.}{fig:cross}
The shipped rule chooses CPU for $E<4,096$ and batches below eight. Counterfactually applied to six H1 cells, it misroutes 3K and 4K, with relative P95 regret of 17.4\% and 23.9\%. For cell $c$ and valid measured routes $\mathcal{R}_c$,
\begin{equation}
R(c)=\frac{T_{\pi(c),95}}{\min_{r\in\mathcal{R}_c}T_{r,95}}-1.
\end{equation}
These compare separate run-level percentiles, not paired queries. Equal weighting of six post-hoc cells is not an application distribution. A profile mechanism in code does not prove adaptive routing reduces regret.
\fig{routing_regret}{H1 RTX 4050, same 100K$\times$384 compact-row cells and FP64 recall 1.0 as Figure~\ref{fig:cross}. Completed-call P95 regret uses median five-run results and frozen 4,096-row rule. Annotations give relative regret. Descriptive counterfactual; no routing traces/confidence bars.}{fig:regret}

\subsection{One-pass and cached eligibility}
H2 uses a separate RTX 4090 with Vulkan, a Ryzen 9 7950X, and driver 570.211.01. Its 100K$\times$384 AG News ablation holds scoring and queries fixed across five runs per mode at $E\in\{1\mathrm{K},4\mathrm{K},6\mathrm{K},100\mathrm{K}\}$. One-pass preparation lowers P95 from 0.131736 to 0.117870 ms at 1K, from 0.305281 to 0.208009 ms at 4K, and from 0.401130 to 0.271618 ms at 6K, reductions of 10.5--32.3\%. The warmed cache lowers the sparse values by 24.3--57.3\%. The dense improvement is smaller: 1.786170 to 1.720476 ms for one-pass preparation, or 1.596024 ms with the cache.
\fig{eligibility_ablation}{H2 \texttt{26e0ccc0}, RTX 4090/Vulkan, 100K$\times$384 AG News, $B=1,k=10$. Median five-run completed-call P95, 25,000 calls/point. Sparse FP64 recall 1.0; dense 0.99998. Sequential modes, logarithmic axes, guide lines, no uncertainty intervals.}{fig:eligibility}
The cache benefit requires an identical predicate and generation; a mutation invalidates that premise. Because H2 has no matched CPU cells and uses another host, it cannot relocate H1's crossover. It does establish that duplicated eligibility work affects a completed call.

\section{Archive-wide routing falsification}
\subsection{Development fit and preregistered failure}
The $EDB$ candidate routes to WGPU when $EDB\geq 10^6$. It was selected from 31 compatible development pairs and classified every development winner. R1 then held hardware and source to its archived RTX 4050/DX12 suite and varied three query shapes, $k\in\{1,10,64\}$, and work on both sides of the frozen threshold. All forced CPU, forced compact-row WGPU, and automatic cells pass their recall and filter gates. The frozen scalar chooses the wrong forced-route winner in 8 of 16 workloads. Median choice regret is only 0.4\% because half the cells have zero regret, but the maximum is 235.7\%; the preregistered 25\% limit rejects the policy.
\begin{table}[H]\centering\small
\caption{Development and held-out evaluation of the frozen $EDB=10^6$ rule. Regret compares separate forced-route median-run P95 values; it does not substitute automatic-route timing.}\label{tab:router-validation}
\begin{tabular}{lrrrrl}\toprule Split & Cells & Wrong & Median \% & Max \% & Decision \\\midrule
\tableinput{tables/final/archive_router_validation}
\bottomrule\end{tabular}\end{table}
\fig{archive_router_failure}{Left: R1, sixteen preregistered RTX 4050/DX12 workloads. The dashed line is the frozen $EDB=10^6$ threshold; ratios below one favor WGPU. Near-equal scalar work produces different winners. Right: separate 100K$\times$384, $B=1,k=10$ H1/P0 cohorts; the P0 source, host, and runtime differ. Lines guide the eye, latencies are not pooled, and no error bars are claimed.}{fig:router-failure}

Near-equal work with different winners alone does not disprove a threshold: a boundary could lie between those values. A genuine ordering contradiction occurs at $k=10$: 768D/$B=16$, $E=61$ favors WGPU at 749,568 units (0.918 versus 0.947 ms), whereas 384D/$B=1$, $E=1,953$ favors CPU at the larger 749,952 units (0.951 versus 1.2404 ms). Thus no increasing $EDB$ threshold perfectly classifies these observed medians. All same-$k$ ordering witnesses depend on the former, only 3.16\%-margin cell; sequential run variability prevents a robust causal conclusion about factorization. A post-hoc threshold of 1,499,904.5 units leaves one error and 3.2\% maximum regret. This is below the original regret limits but fitted to the test set, not new validation. The frozen policy failed; not every scalar policy has been rejected at practical regret limits.

The realized automatic path is also not interchangeable with the forced route selected by a policy. Across R1, automatic P95 ranges from 0.412 to 5.393 times the faster forced-route P95. This spread cannot be labeled routing overhead: sequential run drift, execution state, and path differences remain competing explanations. A router must therefore be evaluated end to end, not only by applying labels to forced-route percentiles.

\subsection{Environment-conditioned crossover and CPU route specialization}
P0 contributes eight CPU/compact-row WGPU pairs on a separate RTX 4090/Linux/Vulkan environment. Its observed reversal lies between 4K and 6K eligible rows, not H1's 2K--3K bracket. At $E=3$K and 4K the two cohorts have opposite winners. Consequently no monotone $E$-only rule can classify all fourteen observed cells. Even the best post-hoc minimax threshold, just above 3K, retains 17.4\% maximum regret; the 4,096 rule retains 23.9\%. This is a falsification over these cells, not attribution to GPU model alone: source, host CPU, OS, API, driver, and execution order also differ.

P2 supplies another conditional route inside the CPU engine. On an A40 host, a certified FP32 candidate stage with FP64 reranking reduces completed-call P95 by 17.0\% at $E=1$K and 28.0\% at 4K, but increases it by 10.0\% at 100K. The measured implementation therefore routes the optimization only through 4,096 eligible rows. Dense single-thread FAISS CPU Flat takes 11.154 ms versus the frozen Qenlo baseline's 15.250 ms, 26.9\% lower latency, while recall is 0.99984 versus 0.99998. This is a scoped library-call comparison, not numerical equivalence or a database ranking.
\begin{table}[H]\centering\small
\caption{P2 A40 selective CPU candidate, 100K$\times$384, $B=1,k=10$, five runs and 5,000 calls/run. Change is optimized/baseline P95 minus one.}\label{tab:phase2}
\begin{tabular}{rrrrr}\toprule $E$ & Baseline ms & Optimized ms & Change \% & Recall@10 \\\midrule
\tableinput{tables/final/phase2_cpu_optimization}
\bottomrule\end{tabular}\end{table}

\section{Predicate representation and batching}
H0 dense GPU endpoints share a recorded GPU revision: shader predicates take 3.2404 ms P95, compact rows 4.2428, and masks 4.4591. At 1K eligible rows, compact rows take 0.6766 ms and shader predicates 2.8832 ms, but the sparse records differ in revision. That difference is context, not a causal representation ablation. The CPU and dense-GPU endpoints also differ. Figure~\ref{fig:representations} preserves these limitations rather than merging H0 with H1.
\fig{windows_strategies}{H0 RTX 4050/DX12, 100K$\times$384 AG News, $B=1,k=10$, five runs/25,000 calls per row, completed-call P95. Dense WGPU forms share \texttt{f1ccc4d}; CPU/sparse records include other revisions. Exhaustive sparse/dense FP64 recall 1.0/0.99998; USearch ANN 0.99224. Descriptive context; no causal cross-revision speedup.}{fig:representations}
Batching affects dispatch amortization and selection; representation affects host preparation and device work. C1--C6 supports both as routing features but changes them together with dimension and selectivity. Device-lab batches use another reporting schema. No same-host factorial study isolates all interactions, and no fixed-$E$, varying-$N$ study establishes eligibility as a statistically superior predictor to size.

The reference C1 and C5 cells sharpen this boundary without isolating it. Both have $E=1$K, $D=128$, $B=1$, and $k=10$, hence equal $EDB$. C1 is a 1K dense collection; C5 obtains the same eligible count through a one-percent predicate over 100K rows. Current CPU P95 rises from 0.0143 to 0.0341 ms, while candidate WGPU remains near 0.094--0.095 ms. Total $N$ and predicate preparation change together, so no causal coefficient follows, but equal eligible scoring work plainly does not imply equal completed-call cost.

\section{ANN and external-engine comparisons}
\subsection{Chroma qualification and historical competitors}
Seven qualified S1 pairs favor candidate WGPU over Chroma HNSW, with P95 ratios from 3.700 to 3233.136. All seven paired Chroma rows have recall 1.0. The largest ratio is the real AG News workload: Chroma takes 555.964 ms versus WGPU's 0.171958 ms. Three attempts fail the gate before held-out timing: C3 and C6 at expansion 128, and deep D2 at 512. Completed USearch D2 has recall 0.3453125 at 83.3133 ms and is unqualified. These rows remain visible.
\fig{small_chroma}{S1 \texttt{730c0500}, seven qualified configuration/workload pairs on A40/RTX 4090, $k=10$; IDs map to matrix/appendix. Logarithmic axis: Chroma/native candidate-WGPU completed-call P95 ratio. Chroma recall printed. Python API, ANN, persistence/filter boundaries differ. Three gate failures excluded from ratios, retained in ledger; no error bars.}{fig:chroma}
H0 Chroma dense reaches 1.1303 ms P95 at recall 0.99414, compared with 3.2404 ms for WGPU predicate evaluation at recall 0.99998. Tuned USearch with expansion 128 takes 3.6245 ms at recall 0.99224. These are operating points with revision and API qualifications, not ANN frontiers. In H3, Linux FAISS HNSW takes 1.4939 ms at recall 0.9994, compared with 12.9014 ms and recall 1.0 for Qenlo CPU. Chroma, USearch, FAISS Flat, Milvus Lite, and LanceDB remain in the appendix. These differences preclude universal rankings.

\subsection{A6000 real and synthetic evidence}
The H4 strict TopK predicate yields $E=10,026$ from 1M 768-dimensional rows. FAISS GPU Flat leads the qualified methods at 0.132222 ms P95, followed by the research CUDA prototype at 0.170722 ms and NumPy at 0.197311 ms. Each has 5,000 calls and FP64 recall 1.0. cuVS's 0.384636 ms row is disqualified because recall is 0.9. The prototype also misses the retained gate requiring a twofold CPU advantage. Eligibility is pre-materialized, and the prototype is Python/PyTorch CUDA rather than shipped Rust/WGPU.

The fresh-seed H5 synthetic cohort differs. FAISS leads at $E=1$K and 10K; the buffered prototype leads at 100K (0.5673 versus 0.7887 ms) and 1M (4.6683 versus 6.3475 ms). The reversal lies between the measured 10K and 100K points. Correctness is unordered top-10 agreement rather than an FP64 oracle. This result is specific to the formulation. Native WGPU has no usable A6000 Vulkan ICD and therefore no native timing in this cohort.
\section{Mutation, reopen, and resource accounting}
The S1 lifecycle uses a durable 10K$\times$384 collection on an RTX 4090 with $B=1$, $k=10$, and no filter. First-search P95 after add-one, delete-one, add-batch, and delete-batch is 0.598, 0.631, 0.740, and 0.819 ms, respectively. Mutation is excluded and has a separate P95 near 22 ms. Three observations of each mutation type report zero rebuilds, and immediate-deletion checks pass. Reopen takes 52.169 ms to open and, separately, 101.296 ms for the first search with one rebuild. Neither value is a combined timer. One observation cannot characterize reopen tail latency.
\fig{small_lifecycle}{S1 reference RTX 4090, durable 10K$\times$384, unfiltered $B=1,k=10$. Bars are first-search latency after mutation/reopen, excluding mutation/opening. Three observations per mutation type give P95; reopen is one observation. Rebuild annotations are fractions, not probabilities. Deletion checks pass. Logarithmic horizontal axis, no error bars.}{fig:lifecycle}
S2 unfiltered WGPU reports 311.181 MB of owned accelerator allocation and a 1.203675 GB process RSS high-water mark. Batch eight reports 314.909 MB and 1.174577 GB, respectively. Units are decimal. RSS includes process and runtime peaks; owned buffers omit driver and runtime allocations. Neither is physical VRAM telemetry. Upload and readback counts are 403,136 and 48 bytes unfiltered, and 517,120 and 65,792 bytes for batch eight. Residency and per-call transfer are separate scopes.
\fig{small_resources}{S2 \texttt{bb195fed}, RTX 4090/Vulkan, generated 100K$\times$768, FP64 recall 1.0. Left: $B=8,k=64,E/N=10\%$, raw phase medians versus call P95. Nested scopes are not additive. Right: two corrected WGPU workloads, zoomed linear latency axis, decimal owned-allocation MB. Different workloads do not establish a memory/latency law or efficiency score. No error bars.}{fig:resources}
Batch-eight median scoring is 0.072 ms, selection 0.613, backend execution 0.826; call P95 is 0.896. Selection exceeds scoring in these diagnostics. That motivates investigation but does not prove another selector improves latency. Capacity, initialization, and warm-state validity remain deployment costs even when arithmetic is inexpensive.

\section{Mobile and heterogeneous devices}
H6 records physical Android: Android 16 (API 36), MediaTek MT6897, Mali-G615 MC6 with Vulkan, app 0.1.0. All 21 aggregate cells report recall@10=1.0 without fallback or failure. CPU wins the 10K quick batch-one cell (5.767 versus WGPU's 12.487 ms P95); WGPU wins the 100K full cell (16.975 versus 29.436) and soak cell (17.678 versus 27.019). Size and suite change together, so this is a qualitative reversal rather than a mobile crossover bracket.
\fig{android_matched}{H6 physical Mali-G615/Vulkan, generated 384D, $k=10$, app 0.1.0; commit unrecorded. Supplied aggregate percentiles: quick/full/soak use 16/64/512 samples/cell. All recalls 1.0; IVF remains approximate. Raw calls and controlled thermal/power records absent. No error bars or current-revision mobile claim.}{fig:android}
At the 100K Android soak point, IVF-Flat takes 21.453 ms P95 and IVF-SQ8 takes 37.090 ms, both slower than exhaustive WGPU at those settings. Dense WGPU batch eight reports 11.131 ms; this batch statistic is not normalized per query. It suggests amortization, not an eightfold per-query gain. The power source is absent, and the thermal state is the uninterpreted value ``0''.

The H7 Intel Arc soak reports CPU at 16.486 ms P95, exhaustive WGPU at 4.444 ms, IVF-Flat at 2.704 ms, and IVF-SQ8 at 22.797 ms for 100K$\times$384. All recalls are 1.0 over 512 samples per cell. Arc adds non-NVIDIA execution evidence and another negative result for universal exhaustive dominance: IVF-Flat wins this cell. One externally labeled ``full'' report contains \texttt{suite=quick}; the JSON is authoritative. Historical Android runs and Kotlin/JVM bindings do not establish current Android packaging. Simulator and macOS builds do not establish physical iOS performance, and no iOS-device or MPS timing is retained.

\section{Negative results and rejected optimization}
Lane-minimum selection retains each lane's minimum and rescans only the winning lane. It reduces repeated scans but changes local state and synchronization. Twelve qualified S0/S1 pairs yield five wins and seven losses, with frozen-baseline to candidate P95 ratios from 0.617 to 1.295. The retained implementation uses the simpler full-rescan selector. Less nominal work did not consistently reduce completed latency. Archives contain other changes too: paired revisions do not prove every difference is solely caused by selection.
\fig{small_selector}{S0 \texttt{f37e744c} versus S1 \texttt{730c0500}, twelve qualified WGPU configuration/workload pairs, FP64-oracle-qualified completed-call P95. Exact workloads/ratios in appendix. Above one favors candidate: five wins, seven losses. RTX 2000 Ada, A40, RTX 4090; no significance test, error bars, or pooled device effect.}{fig:selector}
Four pre-fix 100K$\times$768 failures remain because a 2 GiB benchmark budget did not override the 512 MiB default in-memory admission limit. The corrected source and result archives have distinct hashes; correction does not retroactively qualify the failed rows. Two malformed workload-label invocations are invalid-harness rows, not slow engine results. Chroma gate failures, unqualified USearch, cuVS recall or runtime failures, and unavailable Vulkan paths remain part of the evidence. A controller reaching completion is insufficient for qualification.

R1 is the negative routing result the earlier archive lacked. The frozen $EDB$ candidate passed development and failed held out; it was reverted before release. P2 similarly rejects unconditional use of a sparse CPU optimization. Together with the five-win/seven-loss WGPU selector, these results show three distinct ways that plausible work reduction can fail at the completed-call boundary.

Final known daily Runpod spend is \$0.9400778694252952. Zero campaign pods remain, with ten creations and ten verified deletions. No campaign volume was retained. This is account and daily billing, not per-query cost or a clean marginal campaign cost; billing lag is possible. No cloud work was rerun for authoring.

\section{Deployment implications}
CPU fits tiny eligible workloads when accelerator fixed costs dominate. WGPU can help dense, batched workloads when residency is valid. Tensor execution may fit applications that already own its runtime, but package footprint and persistence-integration costs were not measured. Before preparation is amortized, callers must account for snapshot generation and mutation state.
An explanatory decomposition is
\begin{equation}
T_{\mathrm{call}}=T_{\mathrm{prepare}}+T_{\mathrm{route}}+T_{\mathrm{execute}}+T_{\mathrm{return}}.
\end{equation}
These conceptual, nonoverlapping boundaries are not sums of overlapping counters. Candidate explanatory variables include $E$, $D$, $B$, $k$, representation, and device; R1 rejects one policy based on their product, but does not isolate each factor's causal or predictive contribution. Preparation also depends on total $N$, predicate shape, reuse, generation, and residency. No coefficients are fitted. The decomposition organizes the observations without asserting a transferable latency law.

The reversals and preparation ablation support the central engineering inference that collection size or GPU availability alone is insufficient for routing. R1 goes further: one specific multiplicative predictor is rejected on held-out data. A candidate policy should preserve the variables rather than collapse them, bind calibration to source/device/API identity, include a confidence-aware fallback, and evaluate realized automatic latency after mutations and cold preparation. These requirements define the next experiment; they are not evidence that a learned or calibrated router will succeed. Cloud results cannot select a mobile default.

\section{Threats to validity}
The archive grew sequentially rather than as a randomized factorial study. Dirty trees limit attribution even when hashes and configurations are retained. H0 is cross-revision context. H1 is a post-hoc study on one laptop, one dimension, one batch size, and one corpus, with independently generated metadata. P0 differs from H1 in several environment and source variables; opposite winners falsify universality but do not identify which variable moved the boundary. H2 is sequential, lacks CPU cells, and cannot recalibrate H1. C1--C6 confound collection size, dimension, batching, and eligibility. S2 uses a rejected selector whose final replacement was not rerun in the corrected cells.

The three-run small campaign, five-run historical measurements, and single-reopen observation do not support narrow intervals or broad significance. R1 is held out relative to the selected rule, but its sixteen designed cells are not a random deployment distribution; its post-hoc alternative threshold must not be promoted. Repeated queries are not independent deployment workloads. Aggregate device data cannot yield raw distributions, and thermal and power controls are incomplete. Prototype agreement is weaker than an FP64 oracle. Reference checks do not prove that all ties or workloads agree.

Adapters differ in filtering, selection, normalization, threading, persistence, and API. NumPy uses full lexicographic sorting rather than an optimal top-$k$ baseline. Corrected FAISS is CPU Flat, whereas the historical A6000 result is GPU Flat. Memory scopes differ and omit peak total accelerator use. Concurrency, energy, device loss, cache churn, adversarial predicates, and crash schedules were not benchmarked. Native sweeps lack ANN frontiers and filter-aware baselines. No fastest-database claim follows.

\section{Related work and novelty boundary}
FAISS describes optimized exhaustive and approximate CPU and GPU search and selection \citep{johnson2019billion,douze2024faiss}. HNSW provides graph-based ANN search \citep{malkov2020hnsw}. DiskANN and SPANN address storage-aware regimes \citep{subramanya2019diskann,chen2021spann}; Filtered-DiskANN and ACORN modify graph access under predicates \citep{gollapudi2023filtered,patel2024acorn}.

Filtered-ANN studies establish that attributes, selectivity, query and data conditions, tuning, and implementation affect method choice \citep{iff2026fanns,li2025attribute,shi2025unified}. System analyses examine filtering plans and approximate and exhaustive alternatives \citep{amanbayev2026systems}; the latter two studies are preprints. The observation that selective scans sometimes win is therefore not novel. Qenlo's narrower contribution is the completed-call CPU and portable-GPU boundary, eligibility compilation with generation-bound state, and resource and lifecycle evidence for small collections. The archive lacks broad FANNS coverage and claims no new ANN algorithm.

\section{Conclusions}
The principal confirmatory result is the rejection of the frozen $EDB$ rule: zero errors on 31 development pairs became 8 errors in 16 held-out workloads, with 235.7\% maximum backend-choice regret. Separately, an $E$-only threshold cannot classify opposite observed winners at the same 3K and 4K eligible counts across two environments. No increasing $EDB$ threshold perfectly classifies R1's medians either, but its ordering contradiction depends on a 3.16\%-margin cell, and a post-hoc scalar threshold stays below the original regret limits. These observations do not prove that every scalar policy is practically inadequate or that query-shape factorization caused the differences.

The positive conclusion is a routing architecture, not a fitted policy. Execution decisions must preserve $E,D,B,k$, total $N$, predicate representation, source/device/runtime identity, and warm or mutation state. Eligibility plans reduce avoidable preparation; sparse-only CPU optimization and GPU selection must themselves be routed; and automatic-path latency must be measured directly because it is not interchangeable with forced-route P95. The archive also establishes scoped engineering opportunities: the dense Qenlo CPU path trails single-thread FAISS Flat in P2, and WGPU selection can exceed scoring in the corrected 768D diagnostics.

No retained result demonstrates that a multidimensional production router generalizes. The next valid claim requires a frozen policy, interleaved forced-route and automatic measurements, disjoint held-out devices and query shapes, cold/warm and post-mutation strata, and explicit regret limits. Until that experiment passes, Qenlo has evidence for conditional execution and calibrated fallbacks, not for an adaptive-router advantage or a portable performance promise.

\section*{Reproducibility statement}
Retained evidence only: no product or benchmark changes and no experimental reruns. The archive-wide script verifies four source archives, regenerates the new reductions, and inventories every discovered tracked or selected archive-internal sample series. Ledgers preserve claims, revisions and hashes, environments, provenance, timing, counts, correctness, memory, qualification, and exclusions. Paper-only scripts regenerate reductions, tables, and PDF/PNG figures without replacing raw evidence or historical figures. Appendix~\ref{sec:artifact} and \path{paper/REPRODUCE.md} describe the audit and clean build.
\bibliographystyle{plainnat}
\bibliography{references}
\clearpage
\appendix
\section{Detailed retained tables}\label{sec:tables}
These tables are generated from retained machine-readable evidence. They preserve aggregation within each cohort, and an empty field is not a zero. The full 182-row campaign CSV and JSON include every failed, unavailable, invalid, and unqualified row. The main table explains the difference between native-call and replay sample units. All latency units below are milliseconds.

\subsection{Common small-collection native rows}
\begin{table}[H]\centering\small
\caption{S1 common cells from the matrix. CPU values on A40 are absent for this source role; no substitute baseline-CPU row is silently inserted. Workload definitions appear in the main text. All displayed rows qualify at FP64 recall 1.0.}
\begin{tabular}{llrr}\toprule Configuration & Cell & CPU P95 & WGPU P95 \\\midrule
\tableinput{tables/final/common}
\bottomrule\end{tabular}\end{table}

\subsection{Frozen-baseline versus candidate selection}
\begin{table}[H]\centering\scriptsize
\caption{All twelve qualified S0/S1 pairs. The ratio is frozen-baseline P95 divided by candidate P95. IDs encode row count (r), dimension (d), batch size (b), $k$ when nondefault, and eligible fraction (f). The compound cell combines metadata clauses.}
\begin{tabular}{llr}\toprule Configuration & Workload & Ratio \\\midrule
\tableinput{tables/final/selector}
\bottomrule\end{tabular}\end{table}

\subsection{Small real-embedding replay}
\begin{table}[H]\centering\small
\caption{S1 AG News on the reference RTX 4090, 100K$\times$384, $B=16,k=10,E/N=1\%$. There are three runs and 939 calls per native or replayed engine. Chroma and USearch are ANN; all other engines are exhaustive. External flat and tensor matrices are prefiltered.}
\begin{tabular}{lrr}\toprule Engine & Completed-call P95 & FP64 recall \\\midrule
\tableinput{tables/final/real_small}
\bottomrule\end{tabular}\end{table}

\subsection{Historical Windows endpoints and matched localization}
\begin{table}[H]\centering\small
\caption{H0: 100K$\times$384, $B=1,k=10$, five runs and 25,000 calls per row. Cross-revision endpoint context: dense CPU \texttt{4ae9f1d}, dense WGPU \texttt{f1ccc4d}; sparse CPU and rows \texttt{78e3947}, sparse predicate \texttt{f1ccc4d}. ANN rows retain their own revisions and APIs.}
\begin{tabular}{rlrr}\toprule Eligible & System & P95 & FP64 recall \\\midrule
\tableinput{tables/final/windows_summary}
\bottomrule\end{tabular}\end{table}
\begin{table}[H]\centering\small
\caption{H1: matched recorded revision \texttt{3e2a4a9}, RTX 4050/DX12, AG News, 100K$\times$384, $B=1,k=10$, five runs and 25,000 calls per backend and cell. Post-hoc localization; median run P95.}
\begin{tabular}{rrrr}\toprule Eligible & CPU P95 & Compact WGPU P95 & FP64 recall \\\midrule
\tableinput{tables/final/native_crossover_summary}
\bottomrule\end{tabular}\end{table}

\subsection{Historical eligibility ablation}
\begin{table}[H]\centering\small
\caption{H2, RTX 4090/Vulkan, 100K$\times$384, $B=1,k=10$. Five runs and 25,000 calls per mode and cell, with modes executed sequentially. Percentiles are medians of run percentiles.}
\begin{tabular}{rlrrrr}\toprule Eligible & Preparation & P50 & P95 & P99 & FP64 recall \\\midrule
\tableinput{tables/final/eligibility_ablation_summary}
\bottomrule\end{tabular}\end{table}

\subsection{Archive-wide routing and CPU reductions}
\begin{table}[H]\centering\scriptsize
\caption{R1 all sixteen held-out workloads for the frozen $EDB=10^6$ rule. Work is in millions; regret compares the selected forced route with the faster forced route. Auto/best compares the separately measured automatic P95 with the faster forced P95 and is not interpreted as overhead.}
\begin{tabular}{lrrrrrr}\toprule Workload & $EDB/10^6$ & $k$ & Winner & Choice & Regret \% & Auto/best \\\midrule
\tableinput{tables/final/alpha5_heldout_router}
\bottomrule\end{tabular}\end{table}

\begin{table}[H]\centering\small
\caption{H1/P0 crossover cells used for the equal-$E$ contradiction. Both cohorts use 100K$\times$384, $B=1,k=10$, five runs and 5,000 calls/run/backend; source, host CPU, GPU, OS/API, and driver differ.}
\begin{tabular}{lrrrr}\toprule Environment & $E$ & CPU P95 & WGPU P95 & Winner \\\midrule
\tableinput{tables/final/archive_crossover}
\bottomrule\end{tabular}\end{table}

P2's complete reduction contains the six matched Qenlo rows summarized in Table~\ref{tab:phase2}, a routed dense reference verification at 14.474 ms, and FAISS CPU Flat at 11.154 ms. The latter two are time-separated diagnostics and do not turn run drift into an optimization effect. Full rows, source roles, recall, and archive identity are in \path{research/data/processed/archive-reanalysis/phase2_cpu_optimization.csv}.

\subsection{E0/E2 partial campaign}
\begin{table}[H]\centering\small
\caption{E0/E2 completed block summaries per engine and cell (completed/planned). $f=E/100{,}000$. The archive has no completion marker. The CPU $B=16$, $E=30$K block~1 destination exists without a summary and counts as missing. No value is imputed.}\label{tab:e0e2-coverage}
\begin{tabular}{lrrrrrrrr}\toprule Exp. & $B$ & $E$ & $f$ & CPU & Rows & Predicate & Mask & Total \\\midrule
\tableinput{tables/final/e0e2_coverage}
\bottomrule\end{tabular}\end{table}
\begin{table}[H]\centering\scriptsize
\caption{E2 median block P50 batch-call latency (ms) with 95\% block-bootstrap intervals; five blocks unless $n$ is shown. An $n=1$ cell has no interval. ``missing'' cells were planned and not measured.}\label{tab:e2-p50}
\begin{tabular}{rrllll}\toprule $B$ & $f$ & CPU & Rows & Predicate & Mask \\\midrule
\tableinput{tables/final/e2_latency_p50}
\bottomrule\end{tabular}\end{table}
\begin{table}[H]\centering\scriptsize
\caption{E2 median block P95 batch-call latency (ms) with 95\% block-bootstrap intervals, as in Table~\ref{tab:e2-p50}. At $B=1$, $f=1$, all four engines report recall@10 of exactly 0.99998; all other cells report 1.}\label{tab:e2-p95}
\begin{tabular}{rrllll}\toprule $B$ & $f$ & CPU & Rows & Predicate & Mask \\\midrule
\tableinput{tables/final/e2_latency_p95}
\bottomrule\end{tabular}\end{table}
\begin{table}[H]\centering\scriptsize
\caption{E2 paired ratios against the dense mask, notation as in Table~\ref{tab:e2-paired}. Omitting the exit-code-101 mask block at $B=1$, $f=0.001$ changes the rows/mask P95 ratio from 0.197 to 0.198.}\label{tab:e2-mask}
\begin{tabular}{rrrllll}\toprule $B$ & $f$ & $n$ & Rows/mask P50 & Rows/mask P95 & Predicate/mask P50 & Predicate/mask P95 \\\midrule
\tableinput{tables/final/e2_rows_vs_mask}
\bottomrule\end{tabular}\end{table}
Per-block values, including within-block run spread, process return code, and start time, are in \path{research/data/processed/runpod-e0-e2-analysis-20260924/block_metrics.csv}. The E0 modified-$z$ rule flags CPU block 6 at P95 (0.756 ms) and gpu-rows blocks 5 and 6 at P50 (0.432 and 0.366 ms) and P95 (0.506 and 0.424 ms); none is removed. Two derived files report different bootstrap lower endpoints for the E0 gpu-rows P95 median, 0.614 ms in \path{e0_noise.csv} and 0.625 ms in \path{latency_summary.csv}, because they use different seeds; Table~\ref{tab:e0} uses the former. Minimum per-call recall@10 across 2,181,693 sample rows is 0.9, while every summary passes its aggregate target. Recorded GPU telemetry spans 210--2,775 MHz SM clock, 23--35\,$^\circ$C, and 12.24--131.26 W, as before/after snapshots only. Driver, CUDA, and host-CPU identities are not recorded in the archive.

\subsection{Historical A6000 distributions}
\begin{table}[H]\centering\small
\caption{H4 strict real TopK, 1M$\times$768, $E=10,026$, $B=1,k=10$, 5,000 calls/system. Shared prefiltered matrix, host completed calls. cuVS fails correctness and is excluded from qualified comparisons. CUDA prototype is unshipped research code.}
\begin{tabular}{lrrrr}\toprule System & P50 & P95 & P99 & FP64 recall \\\midrule
\tableinput{tables/final/a6000_strict_summary}
\bottomrule\end{tabular}\end{table}
\begin{table}[H]\centering\small
\caption{H5 generated normalized FP32, 1M$\times$768, $B=1,k=10$, 500 calls/system/cell across five runs. Completed-call distribution percentiles; unordered top-10 agreement, not independent FP64 truth.}
\begin{tabular}{rlrrr}\toprule Eligible & System & P50 & P95 & P99 \\\midrule
\tableinput{tables/final/a6000_exact_summary}
\bottomrule\end{tabular}\end{table}

\subsection{Historical Linux libraries}
\fig{linux_library_context}{H3 \texttt{bb51987}, Linux CPU library context, separately prepared AG News 100K$\times$384, unfiltered $B=1,k=10$. Three runs/1,500 calls per system; completed library-call P95. Recall accompanies latency; algorithms, persistence, and APIs differ. This figure is descriptive within the host, not a product leaderboard.}{fig:linux}
\begin{table}[H]\centering\small
\caption{H3 complete retained library context. ANN recall is part of each operating point; exhaustive engines use different APIs. Native GPU failed Vulkan initialization and has no timed row.}
\begin{tabular}{llrr}\toprule System & Search & P95 & FP64 recall \\\midrule
\tableinput{tables/final/linux_library_summary}
\bottomrule\end{tabular}\end{table}

\subsection{Android aggregate records}
\begin{table}[H]\centering\scriptsize
\caption{H6 all 21 supplied Android cells. $D=384,k=10$; $B$ is batch size and $n$ is the reported sample count. All recall@10 values are 1.0, and all fallback counts are zero. Automatic rows use selective predicates and are not comparable with dense rows. No raw-call reconstruction is possible.}
\begin{tabular}{llrrrrrr}\toprule Suite & Cell & Rows & B & n & P50 & P95 & P99 \\\midrule
\tableinput{tables/final/android_device_lab}
\bottomrule\end{tabular}\end{table}


\subsection{Intel Arc aggregate records}
\begin{table}[H]\centering\small
\caption{H7 all 21 Intel Arc/Vulkan cells, app 0.1.0, $D=384,k=10$. All reported recall is 1.0, with no fallback. The archive contains two quick reports and one soak; labels follow the embedded JSON. Automatic cells are selective and separate from dense rows. Raw series are unavailable.}
\begin{tabular}{llrrrrrr}\toprule Suite & Cell & Rows & B & n & P50 & P95 & P99 \\ \midrule
\tableinput{tables/final/intel_arc}
\bottomrule\end{tabular}\end{table}

\section{Claim ledger and provenance}\label{sec:artifact}
The definitive machine-readable ledger is \path{paper/tables/claim-to-artifact.json}, with a CSV index at the same stem. It incorporates audited historical, campaign, and system claims together with every matrix row. Each measurement records its source, revision or hash, environment, data provenance, workload, filter, timing, sample or repetition unit, oracle, memory scope, qualification, and limitations. Unknown provenance is explicit rather than inferred from a nearby cohort. Full machine paths remain outside the narrow printed tables for legibility.

\begin{table}[H]\centering\small
\caption{Headline claim families and authoritative artifacts. Paths are relative to the repository; complete hashes and per-row fields are in the ledger.}
\begin{tabular}{L{.17\linewidth}L{.37\linewidth}L{.35\linewidth}}
\toprule Claim & Result / qualification & Artifact family \\ \midrule
H1/H1-regret & 2K--3K reversal; up to 23.9\% counterfactual regret, one historical laptop & \path{research/data/raw/2026-09-02-native-crossover/}; per-cell samples/runs/configurations \\
H0 & Historical representation/endpoints; different revisions explicitly separated & \path{benchmarks/2026-08-28/real/} \\
H2 & One-pass/cache preparation; five sequential runs, no CPU recalibration & \path{research/artifacts/qenlo-phase1-eligibility-2026-09-04.tar.gz} \\
S2 & Twelve corrected 768D rows, recall 1.0; candidate selector & Campaign \path{deep768/rtx4090-eu-ro-secure-deep768-admission-fix/} \\
S1-selector & Five wins/seven losses; rejected unconditional optimization & Matrix paired \texttt{baseline-gpu}/\texttt{current-gpu}; source archives \\
S1-real/ANN & Real fixture; seven qualified Chroma pairs; failed gates and USearch visible & Campaign \path{deep/} and \path{reference/}; matrix and raw samples \\
S1-lifecycle & Search-only P95; separate mutation/opening; zero/final-one rebuilds & Reference \texttt{lifecycle-current-gpu} samples/summary \\
S2-resources & Device allocations, process RSS, transfers, nested phases & Corrected native samples, process time logs, matrix \\
\bottomrule\end{tabular}\end{table}
\begin{table}[H]\centering\small
\caption{Historical, platform, accounting, and semantic claim families (continued).}
\begin{tabular}{L{.17\linewidth}L{.37\linewidth}L{.35\linewidth}}
\toprule Claim & Result / qualification & Artifact family \\ \midrule
H4/H5 & Real strict negative; synthetic formulation reversal & \path{benchmark-results/2026-09-02-runpod-a6000-strict-research-gate/}; \path{research/data/raw/2026-09-02-a6000-cuda/heldout/} \\
H3 & Seven CPU library rows, no WGPU timing & \path{benchmark-results/2026-09-02-runpod-archive/retained-archive/artifacts/} \\
H6/H7 & Physical Android and Arc, aggregate-only & \path{research/data/raw/2026-09-02-android/}; \path{benchmarks/2026-08-31/device-lab/intel-arc/reports.json} \\
R1 & Frozen $EDB$ rule: 8/16 wrong, 235.7\% max regret, rejected & \path{research/data/raw/alpha5-router-heldout-rtx4050.tar.gz}; archive reanalysis \\
P0 & Separate 4K--6K reversal; equal-$E$ contradiction with H1 & \path{research/artifacts/qenlo-phase0-runpod-complete-2026-09-04.tar.gz} \\
P2 & Sparse CPU gains, dense regression, FAISS CPU context & \path{research/artifacts/qenlo-phase2-a40-2026-09-04-final.tar.gz} \\
E0 & Same-cell block ratios near $2\times$; GPU bimodality associated with clock snapshots & \path{research/data/raw/runpod-e0-e2-partial-20260924.tar.gz}; \path{research/data/processed/runpod-e0-e2-analysis-20260924/e0_noise.csv} \\
E2 & 227/260 summaries; sparse rows 3.2--6.1$\times$ faster than predicate/mask; no $f^*(16)$ & Same archive; \path{paired_comparisons.csv}, \path{latency_summary.csv}, \path{coverage.csv} \\
E0/E2 witness & Same-host $EDB$ ordering contradiction, post hoc & \path{research/data/processed/runpod-e0-e2-mechanisms-20260924/edb_ordering_witness.csv} \\
Harness & Batch transfer counters overstated $B$-fold; S2 batch-8 corrected & \path{crates/qenlo-bench/src/main.rs}; mechanisms \path{diagnostics_by_cell.csv} \\
Accounting & Account daily spend, zero retained campaign resources & Campaign \texttt{ledger.jsonl}, billing captures, \texttt{campaign-summary.json} \\
Semantics & Current source contracts; not retrospective performance & \path{paper/audit/system-ledger.json} \\
\bottomrule\end{tabular}\end{table}


\subsection{Source and result archives}
The corrected replay records Python 3.12.3, FAISS 1.15.0, NumPy 2.3.3, and PyTorch 2.9.1+cu128. Driver and runtime version labels describe different layers: the NVIDIA driver reports CUDA 13.0, while the PyTorch build targets CUDA 12.8. Per-configuration environment captures remain authoritative for other hosts.

The campaign root is \path{research/artifacts/runpod-small-2026-09-05/}. Full identities are:
\begin{description}
\item[S0 frozen source.] \path{f37e744c7ee4054a74c4b4181f2dde61dade4d677bbc5493b9a36b61d165a45d}
\item[S1 experimental source.] \path{730c050065e6786972a7be4f9bcb619168becdc943ce1122d1583c61f66409e3}
\item[S2 corrected source.] \path{bb195fedb6519c26598c6c103e7e182982c2eddd0a06d9c2cf54eb5cb5ceeb19}
\item[S2 result archive.] \path{9e322d615c0bbed44616faa831493ca030dd49bbce7fe0c4386f49e3b8532fc0}
\item[H2 eligibility archive.] \path{26e0ccc057c59c0fb4687f8d85a923e88b7bb25257e4e551e22cffe13c1b821f}
\item[R1 held-out router archive.] \path{8e3b5fe6ec6ab9f295752ec894c89b094c4ddf4dd0ee38849d9d78c9a7fcc01e}
\item[P0 complete archive.] \path{7292fd1cd0e7669c9febeaa2d49baedee747811ed65f6bc704ecbcb6e61f51ea}
\item[P2 CPU archive.] \path{ac3afda446f4943eab128cff458402f7559be53b96c3cf8ce44fcbbd7f8be4df}
\item[E0/E2 partial archive.] \path{e4ca1336757cf55a5f5150a8bcc601ab3538cb7e9e9c2270c137e5a72fd51ed0}
\item[E0/E2 attested source.] \path{e5d85adbb638caa05a99fb8d1ef0a96147091cf0}
\item[E0/E2 attested binary.] \path{b28e16ffda779304f10bc6fea61ec6fe368ce0923590593e2ab27443bc70285b}
\item[E0/E2 attested dataset.] \path{ba87a322b6846ce225ce54140cf00ac33250271447ef8f9f83df246131719cba}
\item[E0/E2 runner (recomputed).] \path{c0ae079bc4ddbac8ad4efef2a4a25c5be8418e9a1276edb5bcd7dda71a24ae56}
\end{description}
S2 result integrity is distinct from S2 source identity. The corrected result manifest has 74 non-self entries. The full campaign verifier checks eight retained result archives and nine checksum manifests, covering 1,053 entries and skipping four self-manifest entries. It also checks source roles and the identity of the local and baseline WGSL. Checksums establish byte integrity, not scientific qualification.

\subsection{Discrepancies resolved during authoring}
The audit reconciled disagreements instead of choosing flattering prose. It records that (i) ``current'' source archives still contain the rejected selector; (ii) the matrix lifecycle scope string is broader than the separate raw timers; (iii) external replay preparation omits eligibility traversal despite its field name; (iv) completion is not identical to qualification; (v) historical endpoint revisions differ; (vi) Intel Arc's embedded quick label overrides an external full label; and (vii) final spend is account and daily accounting. Bibliographic corrections include the DTU dataset's publication year, the SPANN title, DiskANN author spelling, and the full Filtered-DiskANN author list. Audit notes retain these resolutions.

\subsection{Excluded evidence and unresolved claims}
No raw evidence or old figure was deleted. The historical inventory hashes every top-level figure file, and regenerated figures live in a separate \path{paper/figures/final/} directory. Older dimension, matched-size, metadata-correlation, and heatmap graphics do not establish unmeasured experiments. Redundant crossover and selectivity plots are superseded by explicitly scoped panels. The old architecture figure is retained, while the paper explains canonical state in equations and prose.

An incomplete historical CPU development run and duplicate cloud mirrors are excluded from the counts. H3 native Vulkan initialization failed, so there is no GPU latency to rank. Historical Qdrant was deliberately outside the retained comparison because its local-mode search boundary differed; no conclusion about its performance follows. Provider-compatible TopK experiments used an inclusive predicate that yielded 10,148 eligible rows and supplied truth, whereas strict H4 used 10,026 rows and an independent oracle. These experiments remain exploratory negative context and are not pooled with H4. Their approximately one-percent prototype did not beat FAISS, and other provider cells had correctness disagreements. cuVS numerical and runtime failures remain failures, not qualified exact results.

The archive does not test a complete multidimensional routing surface, controlled fixed-$E$ with varying-$N$, controlled metadata and vector correlation, a tuned filtered-ANN frontier, a successful calibrated router on held-out data, concurrency, energy, total VRAM peaks, long-running mutation churn, arbitrary crash recovery, MPS, or physical iOS. R1 supplies a held-out rejection, not a positive adaptive result. Historical mobile results cannot certify the final revision. The S2 rerun with the final local selector is absent. E2 lacks its batch-16 cells at $E=100$K and most of $E=30$K, an automatic-path arm, continuous GPU telemetry, and repeated noise controls outside one sparse batch-one cell. These gaps constrain the claims rather than licensing extrapolation.

\subsection{Reproduction and audit procedure}
The archive-wide entry point is \path{research/scripts/analyze_full_archive.py}; it verifies R1/P0/H2/P2 hashes, parses archives without extracting over retained data, emits the routing/CPU reductions, and writes a disposition ledger for tracked or selected archive-internal sample series. Other paper-only entry points are \path{paper/scripts/reduce_final_evidence.py}, \path{generate_final_tables.py}, \path{generate_final_figures.py}, and \path{verify_campaign_claims.py}. Historical raw checks are in \path{paper/audit/verify_historical_raw.py}. Reductions go to \path{paper/audit/reduced/} and \path{research/data/processed/archive-reanalysis/}. Android values are regenerated from retained aggregate records rather than unavailable raw series. Numerical tables and figure inputs are traceable through \path{paper/audit/figure-sources.json}. A clean LaTeX build, citation, reference, and box checks, extracted text, and page renders are recorded in the paper verification output. E0/E2 has three read-only stages: \path{research/data/processed/runpod-e0-e2-analysis-20260924/analyze.py} and \path{render_report.py} produce the latency statistics from a Python-extracted copy of the archive; \path{research/scripts/audit_e0_e2_mechanisms.py} reads the tarball in place and produces diagnostics, clock associations, and the ordering witness; and \path{paper/scripts/generate_e0e2_tables.py} writes the table bodies. Cargo tests separately verify the post-audit batch-counter, range-materialization, and required-predicate changes; they do not regenerate or revise campaign measurements. No command in this authoring workflow creates cloud resources or runs a performance benchmark.

\end{document}
